% Part: turing-machines
% Chapter: undecidability
% Section: universal-tm

\documentclass[../../../include/open-logic-section]{subfiles}

\begin{document}

\olfileid{tur}{und}{uni}
\olsection{Mesin Turing Universal}

Ing \olref[enu]{sec} kita wis ngrembug carane
saben mesin Turing bisa diterangake nganggo runtutan
winates wilangan wutuh. Runtutan iku ngode kahanan,
alfabet, kahanan wiwitan, lan instruksi mesin Turing.
Kita uga wis nuduhake yen himpunan kabeh andharan
iki !!{enumerable}. Amarga himpunan andharan mau
!!{denumerable}, ana fungsi !!a{surjective}
saka~$\Nat$ menyang himpunan andharan iku. Fungsi
!!a{surjective} kaya mangkono bisa diasilake,
umpamane, nganggo cara zig-zag Cantor. Kanthi iki,
kita bisa ngenumerasi kabeh (andharan) mesin Turing.
Yen salah siji enumerasi iku wis ditetepake,
kita bisa nyebut mesin Turing kapangka-$1$,
kapangka-$2$, \dots, kapangka-$e$. Wilangan
iki diarani \emph{indeks}.

\begin{defn}
Yen $M$~yaiku mesin Turing kapangka-$e$
(ing enumerasi kang wis ditetepake), kita nyebut
$e$~minangka \emph{indeks} saka~$M$. Kita nulis
$M_e$ kanggo mesin Turing kapangka-$e$.
\end{defn}

Siji mesin bisa nduwé luwih saka siji indeks.
Contone, rong andharan saka~$M$ bisa béda
ing urutan dhaptar instruksine, lan andharan
kang béda iku nduwé indeks kang béda.

Sing wigati, enumerasi andharan mesin Turing bisa
ditetepake saengga andharan~$M$ bisa diitung
kanthi efektif saka indekse, lan sawijining indeks
mesin~$M$ bisa diitung kanthi efektif saka
andharane. Miturut tesis Church--Turing, banjur
ana mesin Turing kang bisa nemokake andharan
mesin Turing kanthi indeks~$e$ lan nulis andharan
iku ing pitane minangka output. Andharan iku
awujud runtutan blok-blok~$\TMstroke$ (kang
nggambarake wilangan wutuh positif ing runtutan
andharan~$M_e$).

Saka kene, lumrah yen kita takon: fungsi apa
saka indeks mesin Turing kang uga bisa diitung
dening mesin Turing? Sipat apa saka indeks mesin
Turing kang bisa diputusake dening mesin Turing?
Tuladhane, fungsi kang nggandhengake indeks~$e$
karo cacahe kahanan mesin Turing kanthi indeks~$e$
bisa diitung dening mesin Turing. Mesin kaya
mangkono, yen diwiwiti ing pita kang ngemot
siji blok $e$~$\TMstroke$, dhisik mbukak kode
$e$ dadi andharane. Saiki andharan iku awujud
runtutan blok-blok~$\TMstroke$ ing pita. Amarga
!!{element} kapisan ing runtutan iki yaiku
cacahe kahanan, mesin mung perlu mbusak kabeh
kejaba blok $\TMstroke$ kapisan, banjur mandheg.

A remarkable result is the following:

\begin{thm}\ollabel{thm:universal-tm} Ana \emph{mesin Turing
  universal}~$U$ kang, yen diwiwiti nganggo input $\tuple{e,n}$ 
  \begin{enumerate}
    \item mandheg yen lan mung yen $M_e$ mandheg nganggo input~$n$, lan
    \item yen $M_e$ mandheg kanthi output $m$, $U$ uga mangkono.
  \end{enumerate}
  Mula $U$ ngitung fungsi $f\colon \Nat \times \Nat \pto \Nat$
  kang $f(e,n) = m$ yen $M_e$, nalika diwiwiti
  nganggo input~$n$, mandheg kanthi output~$m$,
  lan ora ditegesi yen ora mangkono.
\end{thm}

\begin{proof}
  Mbangun~$U$ kanthi nyata iku praktis ora mungkin,
  amarga mesin iki rumit banget. Nanging kita bisa
  njlentrehake garis gedhé carane makarya, banjur
  nggunakake tesis Church--Turing. Nalika diwiwiti,
  pita~$U$ ngemot blok $e$~$\TMstroke$ banjur blok
  $n$~$\TMstroke$. Dhisik mesin ``mbukak kode''
  indeks~$e$ ing sisih tengen input~$n$. Iki
  ngasilake dhaptar wilangan (yaiku blok-blok
  $\TMstroke$ kang dipisahake dening~$\TMblank$)
  kang njlentrehake instruksi mesin~$M_e$. Banjur
  $U$ nulis nomer kahanan wiwitan~$M_e$ lan
  wilangan~$1$ ing pita sisih tengen andharan~$M_e$.
  (Iki uga digambarake kanthi siji simbol bola-bali,
  minangka blok-blok $\TMstroke$.) Sabanjure,
  mesin nyalin input (blok $n$~$\TMstroke$) menyang
  sisih tengen---nanging saben $\TMstroke$ diganti
  blok telung $\TMstroke$ (elinga, nomer simbol
  $\TMstroke$ yaiku~$3$, $1$ yaiku nomer~$\TMendtape$,
  lan $2$ yaiku nomer~$\TMblank$). Ing pucuk kiwa
  runtutan blok-blok iki (kang dipisahake dening
  simbol $\TMblank$ ing pita~$U$), mesin nulis
  siji~$\TMstroke$, yaiku kode kanggo~$\TMendtape$.

  Saiki pita~$U$ ngemot: indeks~$e$, wilangan~$n$,
  nomer kode kahanan wiwitan (yaiku ``kahanan saiki''),
  nomer panggonan wiwitan sirah maca-nulis~$1$
  (yaiku ``panggonan sirah saiki''), lan isi
  wiwitan ``pita'' (runtutan blok-blok~$\TMstroke$
  kang nggambarake nomer kode simbol-simbol~$M_e$
  ---yaiku ``simbol-simbol''---lan dipisahake
  dening~$\TMblank$).

  Banjur mesin iku niru apa kang bakal ditindakake
  $M_e$ yen diwiwiti nganggo input~$n$, kanthi
  langkah-langkah iki:
  \begin{enumerate}
    \item Temokake wilangan~$k$ kanggo ``panggonan
    sirah saiki'' (ing wiwitan yaiku~$1$),
    \item Pindhaha menyang blok kapangka-$k$ ing
    ``pita'' kanggo ndeleng ``simbol'' ing kono,
    \item\ollabel{find-inst}%
    Temokake instruksi kang cocog karo ``kahanan''
    lan ``simbol'' saiki,
    \item Bali menyang blok kapangka-$k$ ing ``pita''
    lan ganti ``simbol'' ing kono nganggo nomer
    kode simbol kang bakal ditulis~$M_e$,
    \item Pindhaha sirah menyang papan kang nyathet
    ``kahanan'' saiki lan ganti nomer ing kono
    nganggo nomer kahanan anyar,
    % OLPL-151: Langkah ngurangi penghitung posisi durung njlentrehake aturan tetep ing kothak 0 yen instruksi mindhah kiwa.
    \item Pindhaha menyang papan kang nyathet
    ``panggonan ing pita'' lan busak
    siji~$\TMstroke$ utawa tambah siji~$\TMstroke$
    (yen instruksine obah kiwa utawa tengen,
    miturut urutane).
    \item Baleni.\footnote{Ing kene kita ngliwati
    sawetara kangelan kang alus. Contone, $U$~bisa
    mbutuhake papan tambahan nalika nambah penghitung
    ``panggonan sirah saiki''---ing kahanan iki,
    kabeh isi ``pita'' kudu dipindhah menyang tengen.}
  \end{enumerate}
  Yen $M_e$ kang diwiwiti nganggo input~$n$
  ora tau mandheg, $U$ uga ora tau mandheg,
  mula outpute ora ditegesi.

  Yen ing langkah~\olref{find-inst} katemu yen
  andharan~$M_e$ ora ngemot instruksi kanggo
  pasangan ``kahanan''/``simbol'' saiki,
  $M_e$ bakal mandheg. Yen iki kedadeyan,
  $U$ mbusak bagean pita ing sisih kiwa
  ``pita.'' Kanggo saben blok telung~$\TMstroke$
  (kang nggambarake siji~$\TMstroke$ ing pita~$M_e$),
  mesin nulis siji $\TMstroke$ ing pucuk kiwa
  pitane dhewe lan mbusak ``pita'' sithik-sithik.
  Yen wis rampung, pita~$U$ ngemot siji blok
  $\TMstroke$ kang dawane~$m$.
  
  % OLPL-152: Rekonstruksi output ing sumber durung netepake carane ngliwati kode panandha pucuk kiwa lan kothak kosong ing pita simulasi sadurunge ngetrapake uji blok telu.
  Yen $U$ nemoni bab saliyane blok telung~$\TMstroke$
  ing ``pita,'' mesin langsung mandheg. Amarga
  ing kahanan iki pita~$U$ ora ngemot siji blok
  $\TMstroke$, outpute dudu wilangan asli,
  yaiku $f(e,n)$ ora ditegesi ing kahanan iki.
\end{proof}

\end{document}
